\originalchapter{动量壳RG与Wilson--Fisher固定点}
\label{ch:wilson}
\source{8.334 L9--L12；采用$u/4!$规范统一图的组合系数。}
\originalsection{快慢模与连接展开}
设$\mathcal H=\int[\frac12(\nabla\vct\phi)^2+\frac r2\vct\phi^2+\frac u{4!}(\vct\phi^2)^2]$，$n$分量，截止$\Lambda$。分$\phi=\phi_<+\phi_>$，慢模$|q|<\Lambda\e^{-\dd l}$，快模在薄壳。有效作用量$\mathcal H_{\rm eff}[\phi_<]=-\log\avg{\e^{-\mathcal H_{\mathrm{int}}}}_>$加慢Gaussian项，展开为$\avg{\mathcal H_{\mathrm{int}}}_>-\frac12\avg{\mathcal H_{\mathrm{int}}^2}_{>,c}+\cdots$。负二阶号来自对数，不能把各阶相同符号堆上去。

Gaussian快模收缩$\avg{\phi_{>,a}(q)\phi_{>,b}(-q)}=\delta_{ab}/(q^2+r)$。一次四腿顶角中收缩两快腿给质量tadpole；两个顶角各留两慢腿并两条快传播子相连给四次顶角修正。只保留连接图避免重复计入真空归一。
\originalsection{一回路系数的计算}
薄壳积分$I_1=K_d\Lambda^{d-2}\dd l/(1+\widetilde r)$、$I_2=K_d\Lambda^{d-4}\dd l/(1+\widetilde r)^2$，$K_d=S_{d-1}/(2\pi)^d$、$\widetilde r=r/\Lambda^2$。为算出$n$的组合系数，写$\phi=\varphi+\chi$（慢、快），四次顶角中恰有两条快腿的部分为
\begin{equation}
V_2=\frac u{24}\int[2\varphi^2\chi^2+4(\varphi\cdot\chi)^2]
=\int\chi_a B_{ab}(\varphi)\chi_b,\qquad
B_{ab}=\frac u{12}(\varphi^2\delta_{ab}+2\varphi_a\varphi_b).
\end{equation}
快场同点收缩$\langle\chi_a\chi_b\rangle=\delta_{ab}I_1$，所以$\langle V_2\rangle=(n+2)uI_1\int\varphi^2/12$；与$r\varphi^2/2$比较即$\delta r=(n+2)uI_1/6$。二阶连接收缩$\langle\chi_a(x)\chi_b(x)\chi_c(y)\chi_d(y)\rangle_c=G_>(x-y)^2(\delta_{ac}\delta_{bd}+\delta_{ad}\delta_{bc})$。两条快线积分给$I_2$，而
\begin{equation}
\operatorname{tr}B^2=\frac{u^2}{144}(n+8)(\varphi^2)^2,\qquad
-\frac12\langle V_2^2\rangle_c=-I_2\int\operatorname{tr}B^2.
\end{equation}
与$u(\varphi^2)^2/24$比较，$\delta u=-(n+8)u^2I_2/6$。这里的$n$是闭合指标环，额外的2和8来自纵向投影与两种跨顶角Wick配对；例如$n=1$得$\delta r=uI_1/2$、$\delta u=-3u^2I_2/2$。零外动量的局部顶角近似要求外部慢波矢远小于薄壳动量。

缩回截止后，以$g=K_du\Lambda^{d-4}$，$\epsilon=4-d$，得到
\begin{align}
\frac{\dd\widetilde r}{\dd l}&=2\widetilde r+\frac{n+2}{6}\frac g{1+\widetilde r}+O(g^2),\\
\frac{\dd g}{\dd l}&=\epsilon g-\frac{n+8}{6}\frac{g^2}{(1+\widetilde r)^2}+O(g^3).
\end{align}
非普适临界质量被tadpole移为负，不能直接把$r=0$当相互作用模型的精确临界温度。固定点$g_*=6\epsilon/(n+8)+O(\epsilon^2)$，$\widetilde r_*=-(n+2)\epsilon/[2(n+8)]+O(\epsilon^2)$。
\originalsection{临界指数与控制参数}
在临界面附近，温度方向本征值$y_t=2-(n+2)\epsilon/(n+8)+O(\epsilon^2)$，所以$\nu=1/2+(n+2)\epsilon/[4(n+8)]+O(\epsilon^2)$。一回路tadpole与外动量无关，故这一阶$\eta=0$。由前章指数关系，
\begin{equation}
\gamma=1+\frac{n+2}{2(n+8)}\epsilon+O(\epsilon^2),\qquad\beta=\frac12-\frac3{2(n+8)}\epsilon+O(\epsilon^2),\qquad\delta=3+\epsilon+O(\epsilon^2).
\end{equation}
这是一种在$\epsilon\ll1$且固定点$g_*=O(\epsilon)$下可控的非Gaussian展开。将$\epsilon=1$代入可估三维，却有未计高阶误差；不应把一回路三维数当精确值。$n=1$给$\nu\simeq7/12$，与精确二维$\nu=1$的远维数结果更不能仅以一阶硬外推。

两回路首次产生外动量依赖的sunset自能。课程给出的标准$O(n)$结果$\eta=(n+2)\epsilon^2/[2(n+8)^2]+O(\epsilon^3)$需同阶波函数与耦合重整化。此处引用其结果定位两回路的作用，主稿完整算到一回路，未把一回路图称作该二阶系数的证明；相应作业中若要求具体二阶壳图，按其模型另算。
\originalsection{其他相互作用与边际维数}
算符$\phi^{2p}$的Gaussian耦合尺度$y_{2p}=d-p(d-2)$，在$d=4$、$p>2$为负；额外梯度也降低工程相关性。因此近四维的普通四次固定点只需有限个相关调参，而不是每个微观项都同等控制临界。但当四次项也被调到0，六次项可变为三临界主项，上临界维数改为3，不能仍删它。

$d=4$时$g'=-(n+8)g^2/6$，$g(l)=g_0/[1+(n+8)g_0l/6]$缓慢趋0。把这个$1/\log$流放入质量与场响应会产生平均场幂乘对数修正，这正是边际项需要非线性流的具体例子。
